\section{Background}
\label{sec:background}

\subsection{Planetary roughness statistics}

Surface roughness has no single value. \citet{shepard2001} compared topographic data from many
natural surfaces and concluded that a roughness measure must be tied to the horizontal scale at
which it is taken, so they proposed reporting scale dependent parameters in a standard form. The
central quantity is the root mean square height difference between points a fixed lag apart. Its
growth with lag is the deviogram, and on a log scale the slope of the deviogram is the Hurst
exponent of a self affine surface. These definitions let workers with different instruments and
sampling report comparable numbers, provided that each number states its lag.

Laser altimetry turned these definitions into maps. \citet{rosenburg2011} computed slopes and
Hurst exponents of the Moon from LOLA over baselines from about 17 m to about 2.7 km. Most of the
surface behaves like a fractal with one or two exponents over that range, and where the exponent
changes, the breakover typically sits near a 1 km baseline \citep{rosenburg2011}. The median Hurst
exponent is 0.95 in the highlands and 0.76 in the maria \citep{rosenburg2011}. \citet{kreslavsky2013}
mapped lunar roughness from LOLA profile curvature at baselines from 115 m to 1.8 km and found that
hectometre and kilometre roughness correlate poorly, because they record different processes and
ages. On Mars, \citet{kreslavsky2000} measured kilometre scale roughness from MOLA as the median
differential slope over baselines from 0.6 km to about 20 km and showed that geological units
have distinct roughness and distinct scale dependence. Field measurements give the same caution
at smaller scales: lava flow roughness behaves as a self affine surface only over a limited range
of lengths, and the smallest scales follow different statistics \citep{campbell1996}. A Hurst
exponent is therefore a description of a stated range of lags, not a property of a surface at all
scales.

\subsection{Radar polarimetry and the circular polarization ratio}

Radar adds information at the scale of its wavelength. Mini-RF on the Lunar Reconnaissance
Orbiter is a hybrid polarity synthetic aperture radar that measures the Stokes parameters of the
echo \citep{nozette2010}. From them, the circular polarization ratio (CPR) of the same sense to
the opposite sense return is formed, and the m-chi decomposition interprets the echo through the degree of polarization and the
Poincar\'e ellipticity parameter \citep{raney2012}. CPR is sensitive to structure near the radar wavelength. Real CPR values above
unity occur on the Moon, on Venus and on terrestrial lava flows, produced by scattering from rock
edges and between rock facets \citep{campbell2012}. Lunar impact melt flows typically have high
CPR, so deposits that look smooth in optical images may be rough at centimetre to decimetre scales
\citep{carter2012}.

CPR is not a direct measure of rock size. Viewing geometry shifts it: across the inner wall of a
bowl shaped crater the local incidence angle changes by about 30 degrees, which can change CPR by
about 0.2 \citep{fa2013}. The same study attributes the enhanced CPR inside anomalous polar
craters mostly to rocks perched on or buried in the regolith, rather than to ice \citep{fa2013}.
Across the lunar highlands, \citet{riveravalentin2026} found that CPR rises quickly with rock
abundance at low abundance, then levels off and can fall over rock rich terrain, so a CPR of 0.8
may correspond to a rock abundance of about 0.7\% or about 10\%. They show that several
polarimetric properties used together resolve this ambiguity \citep{riveravalentin2026}. Mini-RF
mosaics also depend on how terrain effects are handled. \citet{fassett2024} orthorectified the
archived monostatic data onto lunar topography and reduced topographic effects with an empirical
correction based on the known effect of incidence angle.
Learned inversions inherit the same geometry dependence: comparing east and west looking polar
mosaics gives a look direction uncertainty of 1.01 in the real part of the inferred dielectric
constant \citep{shukla2024}.

\subsection{Quantitative analog studies}

Earth surfaces have long served as candidates for interpreting planetary radar and terrain
signatures, but quantitative studies report overlapping populations rather than clean matches. At
the 2014 to 2015 Holuhraun flow field in Iceland, \citet{tolometti2022holuhraun} tested whether
three mapped lava facies differ in UAVSAR CPR and Sentinel-1 cross polarization. The facies were
statistically distinguishable, yet their CPR populations overlap by 88\% to 89\%, so radar can
suggest a facies in probability but cannot map it \citep{tolometti2022holuhraun}. The same study
found only small differences in lidar Hurst exponents between most lava types. On Mars, equatorial
lava flows are radar bright with CPR above one, like blocky terrestrial flows, but terrain models
of 35 flows from HiRISE show that at metre scale they are smoother than blocky flows and closer to
pahoehoe and rubbly flows \citep{rodriguezsanchezvahamonde2021}. Decimetre and metre roughness can
thus point in different directions on one surface. Learned similarity has also been tried:
\citet{meng2024} trained a convolutional network with a discriminative loss to bring Earth and
Mars image patches of the same location together, and used it to select analog areas in the
evaporite region of the Qaidam Basin. These studies treat an Earth surface as a candidate for
interpreting a signature, a use that field work can then test.

\subsection{Self-supervised and domain-adversarial learning}

Joint embedding methods learn representations by making two augmented views of one input agree.
Their failure mode is collapse to a constant output. Contrastive methods such as SimCLR learn by
contrasting augmented views \citep{chen2020}. Barlow Twins drives the cross
correlation of the two views' embeddings towards the identity \citep{zbontar2021}. VICReg adds an
explicit variance term on every embedding dimension to an invariance term and a covariance term
that decorrelates dimensions \citep{bardes2021}. DINO trains a Vision
Transformer by self distillation without labels, and its features work well as nearest neighbour
classifiers \citep{caron2021}. Domain-adversarial training addresses a
different problem: a gradient reversal layer between a feature extractor and a domain classifier
makes the features useful for the main task while uninformative about the domain
\citep{ganin2015}. The two ideas combine naturally when the domains are planetary bodies and the
aim is an embedding of terrain that does not encode which body a window came from.

\subsection{Learned planetary and terrain representations}

The lunar representation models below are trained on lunar data. LunarFM is a multimodal
masked autoencoder that maps 18 channels from six instruments on three lunar missions into a
shared 768 dimensional embedding, with a training set spanning 70 degrees south to 70 degrees
north \citep{gironamata2026}. SomBench assembles more than 30 co-registered lunar layers into
machine learning ready tiles with leakage safe splits \citep{patil2026}, and a lunar foundation
model was pretrained on its tiles at 1 m and 100 m per pixel \citep{fraccaro2026}. Crater retrieval benchmarks show that self-supervised Vision Transformers
with in-domain pretraining outperform larger generic models \citep{fang2026}. On Earth, topo2vec
learns embeddings of elevation data with self-supervision built on the fractal assumption that
structures recur across scales \citep{kavitzky2021}. GEAR retrieves terrestrial topographic analogs
of the Mariana Trench on the Qinghai-Tibet Plateau with a Siamese graph network and an expert
annotated similarity dataset \citep{liu2026}. Optical surveys at the boulder scale complement
these models: a convolutional network mapped about 94 million boulders in images between
60 degrees south and 60 degrees north \citep{aussel2025}. Section~\ref{sec:embedding} compares the
inputs, training signals and invariances of these models with those of the encoder studied here.
